\section{Consequences for theorem design}
\label{sec:discussion}

The calculus changes the unit of progress.  A result is not evaluated only by
the strength of its inequality, but by the number and quality of physical
interfaces it closes.

First, theorem statements should expose their quantifier prefix.  In a
regulated field theory, ``uniform'' must say whether the constant is uniform
in volume, cutoff, boundary state, observable support, and RG history, and in
which order.  The array in Counterexample~\ref{ex:limits} shows why a prose
adjective is insufficient.

Second, retained states should be designed around recovery obligations.  If a
one-block marginal forgets boundary frustration, the correct response may be
to retain a center-flux or contour variable.  If intersection moments forget
adaptive history, one retains a filtration state.  If action charge forgets
normalization, one transports a free-energy or partition-ratio object.  The
state is adequate when the target property becomes either manifest or
recoverable, not when it is merely small.

Third, observable completeness belongs near the beginning of a spectral
programme.  It is inefficient to prove a sharp decay theorem on a core whose
relation to the physical Hilbert space is unknown.  A sectorial result can be
valuable, but it should be named as sectorial.  The hidden-sector
counterexample makes the difference exact.

Fourth, proof discovery and proof certification can use different
representations.  A gauge-fixed chart, numerical experiment, or effective
action may reveal the right invariant.  The final proof may recover it through
a gauge-invariant algebra, an exact disintegration, or a reconstruction
theorem.  The calculus legitimizes this division of labor while recording the
bridge that remains.

Finally, the framework suggests a comparative measure of research maturity.
A programme advances when an arrow moves from heuristic correspondence to
quantified approximation, from approximation to theorem-backed recovery, or
from pointwise recovery to a uniform domain adequate for the physical limit.
That measure is not a scalar score; it is a partial order of closed
obligations.  In Yang--Mills, the present frontier remains the generation of a
full-measure uniform physical rigidity estimate and the construction of the
continuum OS-positive theory.  The representation calculus does not shorten
those analyses, but it prevents progress on a neighboring predicate from
being mistaken for their completion.
